% Generated by code/interviews-1/03-speed-drills/python/gen_drill.py. Do not edit. French copy by draft/pilot-quant-fr/ivgen_fr.py.
\ivfamily{Produits}{6, 10, 17, 18, 21, 28}{Arrondir un facteur vers le haut et l'autre vers le bas, de la même proportion environ~: $38.7 \times 41.2$ est proche de $40 \times 40$.}
\ivfamily{Quotients}{2, 3, 5, 14, 20, 22}{Arrondir le diviseur à un nombre facile, puis corriger du pourcentage dont on l'a déplacé.}
\ivfamily{Racines carrées}{1, 8, 11, 19, 29}{Écrire le nombre comme un carré multiplié par un petit facteur, et déplacer la racine de la moitié de l'écart en pourcentage~: $\sqrt{5300}$ vaut environ $70 + \tfrac{400}{140}$.}
\ivfamily{Intérêts composés}{7, 12, 15, 16, 30}{Prendre $1.05^{10} \approx e^{0.5}$, et retrancher un peu, car le logarithme de $1 + r$ est inférieur à $r$.}
\ivfamily{Exponentielles et logarithmes}{13, 23, 24, 27}{Utiliser $e^x \approx 1 + x + \tfrac{x^2}{2}$ pour un petit $x$. Construire chaque logarithme à partir de $\ln 2$ et $\ln 10$.}
\ivfamily{Pourcentages}{4, 9, 25, 26}{Arrondir le nombre, et non le pourcentage, puis corriger du même petit facteur.}
